%% ============================================================
%%  Section: Descriptive Statistics
%%  Depositor Discipline in Modern Banking (2026-04-06 draft)
%%
%%  References two tables and two figures defined in tables_figures.tex:
%%    tab:descriptive_sumstats_cecl  (tab01_sumstats_20260403)
%%    tab:cecl_predictors_cross_section  (tab02_cecl_predictors_20260403)
%%    fig:cecl_equity_distribution   (fig01_cecl_dist_20260403)
%%    fig:cecl_adoption_quarter_distribution  (fig02_adoption_dist_20260403)
%%
%%  NOTE: The prior desc-stats file used \label{sec:data} which conflicts
%%  with the data section. This file uses \label{sec:descriptive}.
%% ============================================================

\section{Descriptive Statistics}
\label{sec:descriptive}

In this section, we describe the cross-sectional distribution of the treatment
variable, document the concentration of adoption timing, and compare pre-adoption
characteristics of high- and low-CECL banks.  

\subsection{Distribution of the CECL Adjustment and Adoption Timing}
\label{sec:descriptive:distribution}

Figure~\ref{fig:cecl_equity_distribution} plots the cross-sectional distribution
of \texttt{CECL Adj}, the Day-One retained-earnings charge scaled by pre-adoption
book equity, for the community banks in the estimation sample.  The
distribution is right-skewed.  The full-sample mean is 0.28 percent of equity and
the median is essentially zero: roughly half of community banks adopted CECL with
loan portfolios whose forward-looking expected credit losses, under the new
standard, required little incremental reserve above what the incurred-loss
standard had already accumulated.  The right tail is substantial, however.  The
90th percentile of the distribution lies at approximately 2.5 percent of equity,
the threshold for our binary High-CECL indicator.  The banks that exceed
this threshold show a mean adjustment of approximately 3 percent of equity,
representing a material reduction in reported book capital from a single
mandatory accounting event.

Figure~\ref{fig:cecl_adoption_quarter_distribution} shows the distribution of
adoption quarters across the eligible six-quarter window (2022Q3--2023Q4).
Adoption is sharply concentrated in 2023Q1: 3,997 of the 4,320 sample banks
(92.5 percent) appear in their first CECL-compliant Call Report in that quarter,
exactly as prescribed for calendar-year non-public filers.  The remaining
323 banks are distributed across the five adjacent quarters, reflecting
non-calendar fiscal years, early voluntary adoption, and minor filing lags.
This near-uniform adoption calendar strengthens the identification: the
information shock reaches the market in a single concentrated quarter, making
it implausible that the timing itself was chosen in response to bank-specific
deposit conditions.  We nonetheless use bank-specific event time throughout,
assigning each bank its own $t^{*}$, so that the small minority with off-cycle
adoption dates contribute symmetric pre- and post-adoption windows to the panel.

\subsection{Characteristics of High- and Low-CECL Banks}
\label{sec:descriptive:balance}

Table~\ref{tab:descriptive_sumstats_cecl} compares the High-CECL banks
to the Low-CECL banks at the adoption cross-section.  Several patterns
are informative about both the economics of the shock and the plausibility of
the identifying assumption.

High-CECL banks enter the adoption quarter with meaningfully weaker loan quality.
Their nonperforming loan (NPL) ratio is higher compared to Low-CECL banks.  Existing allowances are also higher (1.55 percent vs.\
1.35 percent of loans), confirming that high-CECL banks had already reserved
more under the incurred-loss standard yet still required a material further charge
when forced to estimate expected losses over the full remaining contractual life
of their portfolios.  High-CECL banks are also less well-capitalized, with an
equity-to-assets ratio of 8.9 percent versus 10.3 percent ($-1.46$ percentage
points, $p < 0.01$), consistent with elevated credit risk and thinner capital
buffers.[@@@ rewrite this paragraph without citing numbers. numbers have changed a bit. look at the new table]

CRE loan share does not differ significantly between the two groups (24.2 percent
vs.\ 25.1 percent), suggesting that variation in CECL exposure is more closely
related to loan-level credit quality than to the industry mix of the portfolio.
C\&I loan share and unrealized securities loss ratios are similarly
indistinguishable, a point we return to in Section~\ref{sec:descriptive:predictors}.

The pre-adoption deposit structure of the two groups is comparable on the
dimensions central to our outcome variables.  Uninsured time deposits as a share
of assets average 5.62 percent for High-CECL banks and 5.28 percent for
Low-CECL banks.  The share of all deposits held in uninsured accounts is
marginally lower for High-CECL banks (37.5 percent vs.\ 39.4 percent),
directionally opposite to what a pre-existing outflow narrative would predict.
The identifying assumption for our difference-in-differences design requires
parallel deposit trends between the two groups, not equal baseline levels; the
event-study pre-period coefficients in Section~\ref{sec:results} show no
systematic divergence in the ten quarters before adoption, consistent with this
assumption. [@@@ don't cite specific numbers. describe patterns]

One observable difference in deposit structure does emerge: High-CECL banks hold
more insured time deposits at baseline (9.1 percent vs.\ 7.3 percent of assets,
$p < 0.01$), reflecting these institutions' greater reliance on retail
interest-bearing funding.  Bank fixed effects absorb this level difference, so
our post-adoption estimates capture within-bank changes.

High-CECL banks earn modestly less before adoption: return on assets averages
0.235 percent per quarter versus 0.284 percent ($-0.049$ percentage points,
$p < 0.01$), and interest expense is somewhat higher (0.264 percent vs.\
0.228 percent of assets, $p < 0.01$).  These differences are small in magnitude
and are absorbed by bank fixed effects, but they confirm that High-CECL banks
are a coherent group for which the CECL disclosure carries the most new
information.
[@@@ don't dvelve too much into small differences. this subsection is too long, shorten. don't cite specific numbers. Only do that if there is a very noticable difference]

\subsection{What Predicts the CECL Adjustment?}
\label{sec:descriptive:predictors}

Table~\ref{tab:cecl_predictors_cross_section} estimates cross-sectional OLS
regressions of \texttt{CECL Adj} on bank characteristics at the adoption
quarter, adding controls in three steps.  The results confirm that the CECL
adjustment is driven by credit-risk fundamentals and not by the deposit
composition or securities portfolio characteristics that might independently
generate depositor responses.

The NPL ratio enters positively and significantly in columns~(2) and~(3)
(coefficient $+0.083$ in column~(2), falling to $+0.064$ in column~(3) after
controlling for the allowance ratio): banks with more problem loans required
larger expected credit loss reserves under CECL's forward-looking methodology.
The allowance ratio enters positively and significantly in column~(3)
($+0.159$, $p < 0.01$): even banks that had accumulated relatively large
incurred-loss reserves still required additional charges when expected losses
were estimated over the full contractual life of their loans.  Log assets
enters positively throughout, indicating that larger community banks tended to
carry larger adjustments, though the economic magnitude is modest.

Three non-results are equally important.  CRE loan share, C\&I loan share, and
unrealized securities loss ratio are all statistically insignificant across all
specifications.  The absence of an unrealized loss effect is especially relevant
for the 2023 setting: concerns about securities portfolio losses drove deposit
outflows from SVB and several other large banks in 2023, but our treatment
variable is empirically unrelated to that channel.  The overall fit is
deliberately low: adjusted $R^2$ rises from 0.21 percent in column~(1) to
1.15 percent in column~(3), consistent with the CECL charge conveying
information not easily recoverable from pre-adoption observables.
[@@@ shorten this section too. communicate the main point and move forward, this is not the focus of the paper]
